\Needspace{20\baselineskip}
\section{Search history, stationary rewriting, and the full UAM}
\label{sec:uam-structure}
This section completes the first structural obligation previously listed in
Appendix~\ref{app:remaining-work}. The argument follows the paper's
Lemma~24, the stationary replay proof of Lemmas~27--28, and Lemma~3.
The proof of Lemma~24 records an operation stack rather than repeating the
paper's case analysis informally: each frame stores an actual earlier search
visit with a chronological execution split. The stack is a proof device and
adds no state or transition to the machine.

The stationary theorem is proved first for simultaneous expansion and then
for ordinary shallow plugging. The latter applies directly to the concrete
substitution relation. Independently, the determinism proof identifies a
canonical reduct by erasing the old focus and relabelling the selected variable
as a search focus. It proves that different decompositions of one source have
isomorphic reducts. Consequently the full UAM is deterministic and refocusing
exactly when its extrinsic behaviour has the corresponding property.
These remain conditions on the behaviour; arbitrary compute rules do not
acquire them automatically.

\Needspace{10\baselineskip}
\subsection{Chronological search visits and the operation stack}


\termdef{search-visit}
\begin{definition}[Search visit]
For a run of total length $n$ ending at coordinate $p$, a visit to vertex $v$
is a pair of lengths $a,b$ with $n=a+b$, an initial run of length $a$ to
$(?,v)$, and a suffix of length $b$ from $(?,v)$ to $p$. Thus this predicate
records an occurrence on the actual run, not merely reachability in another run.
\end{definition}
\lean{Graph/SearchHistory}{10}{13}{Linear.SearchHistory.Visit}


\termdef{operation-stack}
\begin{definition}[Operation stack]
For a predicate $P$ on vertices, a stack ends either at the unique input or at
an argument of an operation. An argument frame carries the parent stack at
the operation's source, evidence of $P$ at that source, and membership of the
current vertex in the ordered argument list. In the history invariant, $P$
is the search-visit predicate for the current execution.
\end{definition}
\begin{lemma}[Stack transport and parent recovery]
A pointwise implication $P\Rightarrow Q$ transports a stack. If the current
vertex is an argument of operation $a$, the stack supplies both the parent
stack and the stored predicate at $a$'s source.
\end{lemma}
\begin{proof}
Transport follows by induction on the stack. For parent recovery, the input
case is impossible because an input cannot be an edge target. In an argument
frame, both the stored operation and $a$ target the current vertex. Target
uniqueness identifies their edges; uniqueness of operation views identifies
their sources. Transport the stored parent and predicate along that equality.
\end{proof}
\lean{Graph/SearchHistory}{16}{19}{Linear.SearchHistory.Stack}
\lean{Graph/SearchHistory}{22}{26}{Linear.SearchHistory.Stack.mono}
\lean{Graph/SearchHistory}{29}{36}{Linear.SearchHistory.Stack.parent}


\begin{lemma}[Extending visits]
Appending one search step preserves every recorded visit, increasing its suffix
length by one. A current search coordinate is a visit with empty suffix.
\end{lemma}
\begin{proof}
Keep the initial prefix and append the new step to the old suffix. The length
equation follows by addition. For a current visit, use the entire run as the
prefix and the zero-step run as suffix.
\end{proof}
\begin{definition}[History invariant]
At coordinate $p$ after $n$ steps, require a stack whose frame predicates are
visits on that run, and a visit to the current vertex itself.
\end{definition}
\lean{Graph/SearchHistory}{39}{42}{Linear.SearchHistory.Visit.step}
\lean{Graph/SearchHistory}{45}{47}{Linear.SearchHistory.Visit.current}
\lean{Graph/SearchHistory}{50}{51}{Linear.SearchHistory.Invariant}


\begin{lemma}[History preservation]
Every search step preserves the history invariant. Hence every finite initial
search execution has that invariant.
\end{lemma}
\begin{proof}
First extend all stored visits by the new step. A variable trigger keeps the
same stack and vertex. Entering an operation with no arguments also keeps the
same vertex. Entering its first argument pushes a frame with the already
recorded source visit; the new search position records itself. On return from
an argument, recover the unique parent frame. Advancing to the next argument
replaces the current frame and records the new search position. Finishing
pops the frame and uses its stored source visit. Induct along an arbitrary
suffix to propagate the invariant, adding execution lengths. At the unique
input the stack is empty and the zero-step visit initializes the induction.
\end{proof}
\lean{Graph/SearchHistory}{54}{79}{Linear.SearchHistory.invariant_step}
\lean{Graph/SearchHistory}{82}{87}{Linear.SearchHistory.invariant_steps}
\lean{Graph/SearchHistory}{90}{94}{Linear.SearchHistory.history}


\begin{theorem}[Lemma 24 in coordinates]
If an initial search run of length $n$ ends at $p$ in a mode other than search,
then there exist $a,b$ with $n=a+b$ and $b>0$, such that the run splits at
$(?,p_2)$ into a prefix of length $a$ and suffix of length $b$.
\end{theorem}
\begin{proof}
Take the current-vertex visit from the history invariant. If its suffix length
were zero, its endpoints would coincide, forcing the final mode to be search.
This contradicts the hypothesis.
\end{proof}
\lean{Graph/SearchHistory}{98}{107}{Linear.SearchHistory.revisit}


\Needspace{10\baselineskip}
\subsection{Actual states, restart, and coordinate reflection}

\termdef{search-restart}
\begin{definition}[Search restart]
Restarting a typed coordinate keeps its vertex and changes only its mode to
search. Restarting an actual state erases its focus, restarts its extracted
coordinate, and realizes that coordinate again. This is the paper's change
of focus mode, up to the already proved insertion--erasure isomorphisms.
\end{definition}
\lean{Graph/SearchHistoryStates}{11}{11}{Linear.TypedPosition.restart}
\lean{Graph/SearchHistoryStates}{60}{62}{Linear.State.restart}


\begin{lemma}[Lifting class executions]
A recursively renamed successor of a realized coordinate has a coordinate on
the same original graph. The same holds for every finite execution, preserving
its exact length.
\end{lemma}
\begin{proof}
For one step, factor recursive search into literal search followed by recursive
endpoint isomorphism. Apply the existing coordinate characterization of literal
search and compose the endpoint isomorphisms. Quotient soundness and exactness
turn that isomorphism into equality of state classes. For a finite run, induct
on its length. Lift its first step, rewrite the intermediate class using the
resulting equality, and lift the suffix. Prepend the lifted coordinate step.
\end{proof}
\lean{Graph/SearchHistoryStates}{14}{25}{Linear.TypedPosition.deepSearch_realize_iff}
\lean{Graph/SearchHistoryStates}{28}{39}{Linear.TypedPosition.class_search_iff}
\lean{Graph/SearchHistoryStates}{42}{55}{Linear.TypedPosition.class_steps_lift}


\begin{lemma}[Restart respects representation]
Restart commutes with recursive state isomorphism and therefore descends to
state classes. Restarting a realized coordinate yields the class of its
coordinate restart.
\end{lemma}
\begin{proof}
The erased isomorphism carries the extracted position to the other extracted
position. Resetting both modes preserves this equality. Apply the theorem that
realization commutes with recursive renaming. This proves well-definedness of
the quotient operation. For an inserted focus, use coordinate recovery and
literal realization transport to identify the restarted class.
\end{proof}
\lean{Graph/SearchHistoryStates}{65}{76}{Linear.State.restart_iso}
\lean{Graph/SearchHistoryStates}{80}{82}{Linear.DeepStateClass.restart}
\lean{Graph/SearchHistoryStates}{88}{101}{Linear.TypedPosition.restart_class}


\begin{theorem}[Lemma 24 for actual rooted states]
A rooted state in non-search mode admits an initial search run through its
restart class, followed by a strictly positive search suffix to its own class.
\end{theorem}
\begin{proof}
Lift its rooted execution to coordinates on its erased graph. The endpoint
isomorphism preserves mode, so apply the coordinate revisit theorem. Realize
both parts of the resulting split. Restart well-definedness identifies the
intermediate class with the restart of the original state.
\end{proof}
\lean{Graph/SearchHistoryStates}{107}{129}{Linear.State.rooted_revisit}


\begin{lemma}[Recovering specified coordinates]
Coordinate renaming composes with graph isomorphisms. An isomorphism between
two realized states induces an underlying-graph isomorphism carrying their
specified coordinates to each other. The root class of a realized state is
the class of initial insertion on its original graph. Coordinate executions
commute with recursive renaming.
\end{lemma}
\begin{proof}
For composition, compare the vertex maps on internal, input, and output ports.
For recovery, compose the inverse insertion-recovery isomorphism, the erased
state isomorphism, and the second insertion-recovery isomorphism. The extracted
coordinate equations and inverse laws give the required coordinate equality.
The root claim follows by transporting the initial input, which ordered
boundary preservation fixes. Induction on an execution transports every step
using search renaming.
\end{proof}
\lean{Graph/RootedCoordinates}{11}{18}{Linear.TypedPosition.deepRename_trans}
\lean{Graph/RootedCoordinates}{22}{36}{Linear.TypedPosition.realization_coordinate_iso}
\lean{Graph/RootedCoordinates}{39}{46}{Linear.TypedPosition.root_class}
\lean{Graph/RootedCoordinates}{49}{55}{Linear.TypedPosition.steps_deepRename}


\begin{theorem}[Exact coordinate characterization of rootedness]
A realized state is rooted if and only if its specified coordinate is reachable
by search from the unique input on its original graph.
\end{theorem}
\begin{proof}
The reverse direction is realization of an initial coordinate run. For the
forward direction, identify the root with initial insertion and lift the class
run. Its endpoint may initially be another coordinate with an isomorphic
realization. Recover the underlying isomorphism carrying that coordinate to
the specified one, and transport the entire run. Its initial input is fixed.
\end{proof}
\lean{Graph/RootedCoordinates}{59}{73}{Linear.TypedPosition.rooted_iff}


\Needspace{10\baselineskip}
\subsection{Stationary replay for simultaneous expansion}

\begin{lemma}[Uniqueness of normal runs]
In a deterministic machine, two executions from the same source to normal
forms have equal lengths and endpoints.
\end{lemma}
\begin{proof}
Induct on the first execution. If it has length zero, normality forces the
other run to be empty. Otherwise the other run cannot be empty; determinism
identifies the first successors, and induction identifies the remaining
lengths and endpoints. This uses normality, not successful finality.
\end{proof}
\begin{lemma}[Rooted coordinate properties]
Initial coordinate runs preserve term typing and operation reachability.
\end{lemma}
\begin{proof}
Both hold at the input. Propagate them together along the run using the already
proved one-step typing and operation-reachability invariants.
\end{proof}
\lean{Graph/StationarySearch}{28}{37}{Linear.StationarySearch.run_properties}


\lean{Graph/StationarySearch}{7}{21}{Machine.Steps.normal_unique}


\termdef{stationary-replay}
\begin{definition}[Stationary replay conditions]
Fix a context, an old filling family, and a designated entry vertex. Every
hole filling is one-way. At each other term-typed hole input, a search focus
has an immediate search step to rewrite mode. The old family preserves and
reflects context search away from boundaries; the new family preserves context
search. These last conditions concern corolla copies and label origins, as
defined earlier, rather than an assumed rootedness result.
\end{definition}
A rewrite coordinate is search-normal: replace its mode by its known rewrite
value and apply the no-rewrite-source search theorem.
\lean{Graph/StationarySearch}{44}{44}{Linear.StationarySearch.OneWayHoles}
\lean{Graph/StationarySearch}{48}{51}{Linear.StationarySearch.OtherInputsTrigger}
\lean{Graph/StationarySearch}{54}{59}{Linear.StationarySearch.rewrite_normal}


\begin{lemma}[No boundary before the designated visit]
Consider a rooted prefix ending at a host coordinate, a positive-length suffix
to the designated search entry, and a further suffix to rewrite mode. The
host coordinate is not a search-entry or backtrack-exit boundary of a hole.
\end{lemma}
\begin{proof}
For entry at the designated vertex, the positive suffix would return to the
same search coordinate. Appending the rewrite suffix then gives two runs from
that coordinate to a normal form with different lengths, contradicting normal
run uniqueness. At another input, the immediate-rewrite condition and
search determinism force the next state to be rewrite-normal, so the run
cannot continue to the designated search coordinate. For a backtrack exit,
the rooted prefix supplies term typing and operation reachability. The one-way
filling obstruction rules out that reachable hole output.
\end{proof}
\lean{Graph/StationarySearch}{64}{95}{Linear.StationarySearch.prefix_away}


\termdef{context-prefix-replay}
\begin{theorem}[Extracting and replaying the context prefix]
The prefix preceding the designated search visit is a run in the unchanged
context. It can therefore be replayed after replacing the old filling family.
\end{theorem}
\begin{proof}
Induct on the length of the suffix to the designated visit, retaining the
initial prefix and the eventual rewrite suffix as induction data. At length
zero, injectivity of the host map identifies the context coordinates. At a
successor, the preceding lemma excludes a boundary. Search reflection extracts
one context step and identifies its image with the actual successor. Extend
the initial prefix by that step and apply induction to the remaining suffix.
For replay, apply Lemma~24 to the initial run ending at the designated rewrite
coordinate. Extract its prefix in the context, then transport that prefix
through the new filling family.
\end{proof}
\lean{Graph/StationarySearch}{98}{121}{Linear.StationarySearch.extract_prefix}
\lean{Graph/StationarySearch}{125}{137}{Linear.StationarySearch.replay}


\begin{theorem}[Rooted stationary endpoints]
The actual search-focused replacement is rooted whenever the old
rewrite-focused state is rooted and the stationary replay conditions hold.
Independent recursive renaming of both endpoints preserves this conclusion.
\end{theorem}
\begin{proof}
Reflect rootedness to the specified old coordinate, replay its prefix, and
realize the resulting initial coordinate run on the replacement. For renamed
endpoints, quotient soundness identifies each endpoint class with its realized
representative; transport the rooted predicate along these equalities.
\end{proof}
\lean{Graph/StationarySearch}{142}{155}{Linear.StationarySearch.rooted}
\lean{Graph/StationarySearch}{159}{183}{Linear.StationarySearch.rooted_iso}


\Needspace{10\baselineskip}
\subsection{Moving an inserted focus through ordinary plugging}

\begin{definition}[Reassociation and focused fragments]
The reassociation bijection sends $A+(B+D)$ to $(A+B)+D$ by preserving each
summand; its inverse reverses this operation, with both inverse laws by cases.
The fragment inserts a focus at a chosen input of a replacement graph. The
whole construction inserts it at the corresponding port after plugging.
\end{definition}
\lean{Graph/PlugFocusInsertion}{9}{13}{Linear.PlugFocusInsertion.reassoc}
\lean{Graph/PlugFocusInsertion}{21}{21}{Linear.PlugFocusInsertion.fragment}
\lean{Graph/PlugFocusInsertion}{22}{25}{Linear.PlugFocusInsertion.whole}


\begin{lemma}[Port identities]
Reassociation carries retained host ports to their focused whole-graph ports,
and inserted old ports to the same whole-graph ports. For a non-output port,
it also commutes with splitting at the selected input.
\end{lemma}
\begin{proof}
The first two claims are cases on the port constructors. For splitting, separate
the selected input from all other ports. The selected input goes to the fresh
singleton on both sides. For another non-output port, injectivity of insertion
on source-side ports excludes equality with the selected whole-graph port;
both splits then reduce to the old-port identity. This restriction avoids
assuming that the complete insertion map is injective for self-looping holes.
\end{proof}
\lean{Graph/PlugFocusInsertion}{27}{30}{Linear.PlugFocusInsertion.host}
\lean{Graph/PlugFocusInsertion}{32}{39}{Linear.PlugFocusInsertion.insert_host}
\lean{Graph/PlugFocusInsertion}{41}{53}{Linear.PlugFocusInsertion.insert_split}


\begin{theorem}[Focus insertion commutes with plugging]
Plugging the input-focused replacement is isomorphic to inserting the focus
at its attached port in the filled whole graph.
\end{theorem}
\begin{proof}
Use reassociation for both the internal-vertex and edge carriers. Labels are
unchanged. A retained context source cannot meet the selected input port:
source uniqueness would identify its edge with the removed hole. Its source
list therefore uses the host identity. Sources of inserted edges use the
split identity; their ports are non-outputs by well-formedness. Targets use
the old-port identities. The new focus's singleton source and target satisfy
the same equations directly.
\end{proof}
\lean{Graph/PlugFocusInsertion}{57}{89}{Linear.PlugFocusInsertion.iso}


\Needspace{10\baselineskip}
\clearpage
\subsection{Reflecting search through one shallow filling}

\begin{definition}[Host coordinates and boundary]
A host coordinate is mapped through the retained-vertex injection. Its boundary
predicate means either search mode at a source of the selected hole, or
backtrack mode at a target of that hole.
\end{definition}
\begin{lemma}[Reading back retained operations]
An operation in the filled graph whose edge is retained has exactly the source
and ordered arguments of an original context operation, mapped by the host
injection.
\end{lemma}
\begin{proof}
Its label supplies an operation view on the original edge. Plug that view back
into the whole graph; it has the same edge as the given view. Uniqueness of
operation data gives the source and argument equations.
\end{proof}
\lean{Graph/PlugSearchReflection}{12}{13}{Linear.PlugSearchReflection.position}
\lean{Graph/PlugSearchReflection}{29}{30}{Linear.PlugSearchReflection.Boundary}
\lean{Graph/PlugSearchReflection}{16}{26}{Linear.PlugSearchReflection.operation_readback}


\begin{lemma}[Inserted incidence meets the boundary]
An inserted source meeting a host vertex meets an input of the replacement;
an inserted target meeting a host vertex meets an output.
\end{lemma}
\begin{proof}
Unpack membership in the mapped incidence list. Source ports cannot be outputs,
and target ports cannot be inputs. The existing insertion/host intersection
lemmas identify the corresponding boundary index, whose attachment belongs to
the hole's original source or target list.
\end{proof}
\lean{Graph/PlugSearchReflection}{33}{37}{Linear.PlugSearchReflection.source_boundary}
\lean{Graph/PlugSearchReflection}{40}{44}{Linear.PlugSearchReflection.target_boundary}


\begin{theorem}[Ordinary search reflection]
Away from the selected hole's boundary, every search step of the filled graph
comes from a context step, with exactly the corresponding successor coordinate.
\end{theorem}
\begin{proof}
Inspect the variable-trigger, operation-entry, and operation-return rules.
For a retained edge, recover its label and singleton source, or read back its
operation view. In the return case, injectivity of the host map recovers the
indexed argument lookup. These data construct a context step. Its plugged
step and the original step have the same successor by search determinism.
An inserted interacting edge instead forces the appropriate boundary case by
the preceding incidence lemmas, contrary to the hypothesis.
\end{proof}
\lean{Graph/PlugSearchReflection}{47}{107}{Linear.PlugSearchReflection.reflect}


\Needspace{10\baselineskip}
\subsection{The stationary criterion for ordinary contexts}

\begin{lemma}[Shallow transport and one-way exclusion]
Every context search run survives filling its hole. A one-way filling rules
out an operation-reachable term vertex at a hole output.
\end{lemma}
\begin{proof}
Transport the run by induction, using preservation of a context search step
under plugging. For exclusion, write the hole target as an indexed output
port, extract its operation path inside the replacement, transfer the term
type through the boundary signature, and apply one-wayness.
\end{proof}
The other-input condition for one hole requires immediate rewrite at every
term-typed source distinct from the designated source.
\lean{Graph/StationaryPlug}{14}{20}{Linear.StationaryPlug.steps}
\lean{Graph/StationaryPlug}{23}{30}{Linear.StationaryPlug.oneWay_target}
\lean{Graph/StationaryPlug}{33}{35}{Linear.StationaryPlug.OtherInputsTrigger}


\begin{theorem}[Ordinary stationary prefix replay]
For a one-way filling satisfying the other-input condition, the prefix leading
to the designated search position before a rewrite is an unchanged-context
run. It replays with any new filling of the same interface.
\end{theorem}
\begin{proof}
Use the three boundary exclusions already described for simultaneous expansion:
normal-run uniqueness at the same input, immediate rewrite at another input,
and one-way exclusion at an output. The single-hole reflection theorem then
extracts each context step by induction on the remaining prefix. Lemma~24
supplies the chosen search visit, and ordinary plugging transports the
extracted run to the new filling. No proof of rootedness of the reduct is
included among these assumptions.
\end{proof}
\lean{Graph/StationaryPlug}{40}{71}{Linear.StationaryPlug.prefix_away}
\lean{Graph/StationaryPlug}{74}{97}{Linear.StationaryPlug.extract_prefix}
\lean{Graph/StationaryPlug}{101}{113}{Linear.StationaryPlug.replay}


\begin{lemma}[Closing a simple context]
A focus-free, hole-free replacement plugged into a focus-free simple context
produces a focus-free, hole-free graph.
\end{lemma}
\begin{proof}
Focus freedom is checked separately on retained and inserted edges. A hole
occurrence either comes from the replacement, contradicting its hole freedom,
or from a retained context edge. Simplicity forces that edge to be the selected
hole, contradicting retention.
\end{proof}
\termdef{stationary-rewrite}
\begin{theorem}[Stationary rewriting preserves rootedness]
Let $C$ be a focus-free simple context and $H$ a focus-free, hole-free one-way
filling. Fix an input $i$ of term type and impose the immediate-rewrite
condition at other term inputs. Replacing $C[\mathsf{rewrite};_i H]$ by
$C[?;_i K]$, for any focus-free, hole-free $K$ of the same interface, preserves
rootedness, allowing arbitrary literal endpoint names.
\end{theorem}
\begin{proof}
The preceding lemma certifies both erased whole graphs. Commutation of focus
insertion with plugging identifies the actual endpoints with their specified
whole-graph coordinates. Reflect the source's rootedness to a coordinate run,
apply ordinary stationary replay, and realize the resulting target run.
Transport both endpoint predicates through their focus-preserving isomorphisms.
\end{proof}
\lean{Graph/StationaryPlug}{117}{129}{Linear.StationaryPlug.filled_conditions}
\lean{Graph/StationaryPlug}{134}{170}{Linear.StationaryPlug.rooted}


\Needspace{10\baselineskip}
\subsection{Rootedness of the concrete substitution rule}

\termdef{blocked-input}
\begin{lemma}[Blocked inputs imply one-wayness]
If no operation edge can leave a term input, the graph is one-way.
\end{lemma}
\begin{proof}
Induct on the length of a reverse operation trace ending at that input. An
empty trace starts at the input. For a nonempty trace, induction identifies
the source of its last relevant edge with the input, contradicting the blocked
operation-source condition. Thus an operation path from that input to a
boundary output would identify an input with an output, which is impossible.
\end{proof}
\begin{corollary}[The variable prefix is one-way]
Prefixing any one-input graph by the unary variable edge produces a one-way graph.
\end{corollary}
\begin{proof}
There is only one input index. The new variable edge owns that input's source
incidence. Any operation leaving it would be the same edge by source
uniqueness, contradicting its variable label.
\end{proof}
\lean{Graph/SubstitutionRefocusing}{11}{24}{Linear.oneWay_of_blocked_inputs}
\lean{Graph/SubstitutionRefocusing}{29}{39}{Linear.Substitution.prefix_oneWay}


\begin{lemma}[Interface and single-input facts]
A one-input hole has no non-designated input, so its other-input condition is
vacuous. Reindexing equal interfaces commutes with initial focus insertion
and preserves one-wayness and hole freedom.
\end{lemma}
\begin{proof}
For the first statement, recover the boundary index of a source occurrence;
the input signature has length one, so that index is the designated one.
For interface transport, eliminate the two interface equalities. The claims
reduce to reflexive reboundary transport; focus insertion uses congruence of
the graph-valued constructor. The same equality elimination proves preservation
of one-wayness and closure.
\end{proof}
\lean{Graph/SubstitutionRefocusing}{42}{51}{Linear.Substitution.single_input}
\lean{Graph/SubstitutionRefocusing}{54}{63}{Linear.Substitution.transport_insert}
\lean{Graph/SubstitutionRefocusing}{66}{68}{Linear.Substitution.oneWay_transport}
\lean{Graph/SubstitutionRefocusing}{71}{73}{Linear.Substitution.holeFree_transport}


\begin{theorem}[Concrete substitution preserves rootedness]
Every transition of the concrete substitution relation preserves rootedness.
\end{theorem}
\begin{proof}
Unpack its simple context, bound graph $G_S$, and endpoint isomorphisms.
Take the old filling to be the variable-prefixed $G_S$, and the new filling
to be $G_S$, both transported to the hole's interface. The old filling is
one-way by the variable-prefix corollary, and the single-input lemma supplies
the other-input clause. Focus freedom and closure follow from the existing
prefix theorems and interface transport. Transported focus insertion identifies
the two original fragments with the stationary fragments. Apply the ordinary
stationary rootedness theorem to the given endpoints.
\end{proof}
\paragraph{Relation to the paper's stability assumption.}
Lemma~3 of the paper assumes that $G_S$ is stable for its refocusing conclusion.
For the unary linear substitution rule formalized here, the argument above is
stronger: the variable prefix itself is one-way, independently of $G_S$, and
there are no other inputs to check. Stability is neither added to the
transition definition nor replaced by an assumed refocusing field. This
observation concerns this rule; it is not a statement about the copying UAM.
\lean{Graph/SubstitutionRefocusing}{80}{104}{Linear.State.substitution_rooted}


\Needspace{10\baselineskip}
\subsection{Determinism despite different context decompositions}

\begin{lemma}[Focus preservation from recursive graph isomorphism]
A recursive graph isomorphism between uniquely focussed graphs preserves their
distinguished focuses and modes.
\end{lemma}
\begin{proof}
Primitive-label preservation carries the first focus label to the image edge.
Uniqueness of the second focus identifies that image edge and its mode.
\end{proof}
\begin{definition}[Resuming a variable]
On a graph with a selected variable edge, change only that edge's label to a
search focus. The old and new signatures are both one term input and one term
output, so the graph remains well-typed.
\end{definition}
\lean{Graph/SubstitutionDeterminism}{9}{13}{Linear.Focussed.deepIsoOfGraph}
\lean{Graph/SubstitutionDeterminism}{17}{27}{Linear.Substitution.resumeVariable}


\begin{lemma}[Resumption respects isomorphism]
An isomorphism carrying one selected variable to another induces an isomorphism
between their resumed graphs.
\end{lemma}
\begin{proof}
Keep the vertex and edge bijections and all incidence equations. At the selected
edge both new labels are the same search focus. At another edge, injectivity
prevents its image being the selected target edge, so the original recursive
label relation still applies.
\end{proof}
\termdef{substitution-presentation}
\begin{definition}[Substitution presentation]
A presentation consists of the erased base graph, focus freedom and closure,
a rewrite-mode coordinate, its adjacent variable edge and singleton source
equation, an isomorphism from realizing that coordinate to the actual source,
and an isomorphism from the variable-resumed base to the actual target.
These are proved data extracted from a substitution transition.
\end{definition}
\lean{Graph/SubstitutionDeterminism}{30}{52}{Linear.Substitution.resumeVariable_iso}
\lean{Graph/SubstitutionDeterminism}{56}{66}{Linear.Substitution.Presentation}


\begin{theorem}[Uniqueness from two presentations]
Two presentations of the same source yield recursively isomorphic focussed targets.
\end{theorem}
\begin{proof}
Compose their source isomorphisms and recover the underlying-graph isomorphism
carrying the two selected coordinates to each other. The first variable's
mapped singleton source is therefore the source of the second variable.
Source uniqueness identifies the two edges. Resumption respects this
isomorphism; compose it with the two target isomorphisms. The resulting graph
isomorphism preserves focus by uniqueness.
\end{proof}
\lean{Graph/SubstitutionDeterminism}{70}{83}{Linear.Substitution.Presentation.deterministic}


\begin{lemma}[Presentations are constructed from substitution]
Relabelling the selected variable in a plugged variable prefix gives a graph
isomorphic to the ordinary reduct. Every substitution transition consequently
supplies a presentation.
\end{lemma}
\begin{proof}
For the first claim use identity bijections: the two graphs have exactly the
same carriers, vertex labels, and incidence lists. Only the selected variable
label changes to a search focus. Inspect retained context edges, old filling
edges, and the added singleton separately. Reboundary transport merely maps
the same ports on both sides. For extraction, use the plugged variable prefix
as base graph, certify its focus freedom and closure, and choose its new
variable edge. Focus-insertion commutation and interface transport give the
source isomorphism. The relabelling lemma gives the target isomorphism.
\end{proof}
\begin{theorem}[Substitution determinism]
Two concrete substitutions from the same state have targets isomorphic at all
box depths, preserving their focuses. Their contexts and fragment interfaces
need not coincide.
\end{theorem}
\begin{proof}
Extract both presentations and apply their uniqueness theorem.
\end{proof}
\lean{Graph/SubstitutionDeterminism}{87}{149}{Linear.Substitution.resume_prefix_iso}
\lean{Graph/SubstitutionDeterminism}{153}{177}{Linear.Substitution.presentation}
\lean{Graph/SubstitutionDeterminism}{183}{188}{Linear.State.substitution_deterministic}


\Needspace{10\baselineskip}
\subsection{The full UAM structural theorem}

\termdef{determinism-modulo}
\begin{definition}[Determinism modulo isomorphism]
A relation on actual named states is deterministic when two successors of the
same source are related by a recursive focus-preserving isomorphism.
Raw carrier names are not required to be equal.
\end{definition}
\termdef{actual-refocusing}
\begin{definition}[Refocusing of actual-state transitions]
A relation is refocusing when it sends every rooted source to a rooted target,
using the rooted predicate already defined on recursive state classes.
\end{definition}
\lean{Graph/UAMStructure}{8}{9}{Linear.DeterministicModulo}
\lean{Graph/UAMStructure}{12}{13}{Linear.RefocusingStates}


\begin{theorem}[Full-UAM determinism]
The full UAM is deterministic modulo recursive isomorphism if and only if its
extrinsic behaviour is deterministic in that sense.
\end{theorem}
\begin{proof}
For sufficiency, inspect the two transitions. Two search steps agree by search
determinism. Two substitutions agree by the new substitution theorem. Two
extrinsic steps agree by the supplied behaviour condition. Search and rewrite
cannot coexist at a source because their focus modes are exclusive. A
substitution and computation cannot coexist because source uniqueness would
identify a variable with an active operation. These exclusions discharge all
mixed cases. Necessity follows because every behaviour step is a UAM step.
\end{proof}
\lean{Graph/UAMStructure}{27}{40}{Linear.UAM.deterministic}
\lean{Graph/UAMStructure}{43}{47}{Linear.UAM.deterministic_iff}


\begin{theorem}[Full-UAM refocusing]
Intrinsic transitions preserve rootedness. The full UAM is refocusing if and
only if its extrinsic behaviour preserves rootedness.
\end{theorem}
\begin{proof}
Search preserves rootedness by appending its step to the rooted run.
Substitution preserves it by the stationary argument. These are the two
intrinsic cases. The full machine adds only extrinsic transitions, handled
by the behaviour hypothesis. Conversely, restrict a full refocusing law to
those extrinsic transitions.
\end{proof}
\lean{Graph/UAMStructure}{19}{23}{Linear.UAM.intrinsic_refocusing}
\lean{Graph/UAMStructure}{50}{54}{Linear.UAM.refocusing}
\lean{Graph/UAMStructure}{57}{61}{Linear.UAM.refocusing_iff}


\begin{corollary}[Rooted execution invariant]
Under the behaviour refocusing condition, every finite full-UAM execution
preserves rootedness. Every execution from an initial state therefore ends
in a rooted state. The empty-behaviour machine is both deterministic and
refocusing without additional assumptions.
\end{corollary}
\begin{proof}
Induct on the full execution and apply one-step refocusing at each successor.
Initial states are rooted by the zero-step initial run. For empty behaviour,
both required behaviour properties follow by elimination of its false step
relation; apply the two structural theorems.
\end{proof}
\lean{Graph/UAMStructure}{64}{68}{Linear.UAM.rooted_execution}
\lean{Graph/UAMStructure}{71}{73}{Linear.UAM.initial_execution_rooted}
\lean{Graph/UAMStructure}{76}{79}{Linear.UAM.empty_structural}


\Needspace{10\baselineskip}
\subsection{Regression executions and the remaining work}

The variable-prefix example has an actual one-step initial search run to its
rewrite-focused substitution source, obtained by the variable trigger rule.
Realization proves that source rooted. The concrete substitution theorem then
proves its reduct rooted, and the determinism theorem identifies the reduct
class of every alternative substitution decomposition. Applying Lemma~24 to
the rooted source produces a strictly positive restart suffix; this exercises
the theorem on an actual execution rather than an empty hypothesis.
\lean{Graph/UAMStructureExamples}{9}{15}{Linear.UAMStructureExamples.source_rooted}
\lean{Graph/UAMStructureExamples}{18}{19}{Linear.UAMStructureExamples.substitution_rooted}
\lean{Graph/UAMStructureExamples}{22}{25}{Linear.UAMStructureExamples.substitution_unique}
\lean{Graph/UAMStructureExamples}{28}{35}{Linear.UAMStructureExamples.source_revisit}


The earlier singleton extrinsic behaviour has a genuine compute transition.
Its two successor witnesses identify the same target, proving determinism.
That target is initial, so it is rooted independently of the source, proving
refocusing. Applying the full-UAM theorems supplies both structural properties
for this nonempty behaviour.
\lean{Graph/UAMStructureExamples}{38}{40}{Linear.UAMStructureExamples.behaviour_deterministic}
\lean{Graph/UAMStructureExamples}{43}{45}{Linear.UAMStructureExamples.behaviour_refocusing}
\lean{Graph/UAMStructureExamples}{48}{50}{Linear.UAMStructureExamples.full_structural}


\paragraph{Completed obligation and honest scope.}
Appendix~\ref{app:remaining-work} now has two items: constructing the complete
concrete specimen interface and its observation, safety, robustness, and lifting
bridges; and instantiating the final contextual robustness theorem. The present
results prove actual transition properties, not the missing specimen laws.
They do not assert that raw named-state successors are literally equal, or
that arbitrary behaviour is deterministic or refocusing. The estimate of
completed mathematical effort is about 88\%, with a plausible range of
74--93\%; this is a subjective planning estimate, not a declaration count.

